Tercera diferencia: con un millón de pilas, en el paper Autores lidera con 129 millones de operaciones por segundo. Aquí, cuatro algoritmos se quedan entre 53 y 57. Un empate.

Cuando cuatro algoritmos tan distintos chocan contra el mismo techo, lo razonable es pensar que el techo no es de los algoritmos, sino de la máquina.

La sospecha: la memoria RAM. Con un millón de pilas, los datos no caben en la caché y casi cada operación va a la RAM. Este PC tiene un solo módulo, un único canal; el paper tenía cuatro chips con su memoria cada uno.

Una sospecha no basta. Lo medimos.
